% ============================================================
%  RE-VALUE --- Deliverable D4
%  Gruppo 21 --- Ingegneria del Software, UniTrento
%  A.A. 2025/2026 
% ============================================================
\documentclass[11pt,a4paper]{article}

\usepackage[utf8]{inputenc}
\usepackage[T1]{fontenc}
\usepackage[italian]{babel}
\usepackage[top=2.5cm,bottom=2.5cm,left=2.5cm,right=2.5cm]{geometry}
\usepackage[hidelinks]{hyperref}
\usepackage{longtable}
\usepackage{booktabs}
\usepackage{tabularx}
\usepackage{array}
\usepackage{pgfplots}
\usepackage{tikz}
\usetikzlibrary{arrows.meta,positioning,shapes.geometric,fit,backgrounds}
\usepackage{xcolor}
\usepackage{microtype}
\usepackage{fancyhdr}
\usepackage{float}
\usepackage{enumitem}
\usepackage{parskip}
\usepackage{seqsplit}

\pgfplotsset{compat=1.18}
\setlength\hfuzz{30pt}
\AtBeginDocument{\setlength\hfuzz{30pt}}

\definecolor{blurevalue}{HTML}{2563EB}
\definecolor{greenrevalue}{HTML}{16A34A}
\definecolor{graylight}{HTML}{F3F4F6}

\pagestyle{fancy}
\fancyhf{}
\lhead{\small RE-VALUE --- Deliverable D4}
\rhead{\small Gruppo 21 --- Ingegneria del Software, UniTrento}
\cfoot{\small\thepage}
\renewcommand{\headrulewidth}{0.4pt}

\newcommand{\code}[1]{\texttt{#1}}
\newcolumntype{L}[1]{>{\raggedright\arraybackslash}p{#1}}
\newcolumntype{C}[1]{>{\centering\arraybackslash}p{#1}}

% ==============================================================
\begin{document}
\hypersetup{pageanchor=false}

\begin{titlepage}
  \centering
  \vspace*{2.5cm}
  {\Huge\bfseries RE-VALUE\par}
  \vspace{0.6em}
  {\large\itshape Piattaforma per il riuso locale di oggetti\par}
  \vspace{1.5em}
  \rule{\linewidth}{0.6pt}\\[0.8em]
  {\LARGE\bfseries Report Milestone \#4 --- Deliverable D4\par}
  \rule{\linewidth}{0.6pt}
  \vspace{1.5em}
  {\large Ingegneria del Software --- A.A.\ 2025/2026\par}
  {\large Università degli Studi di Trento\par}
  \vspace{2em}

  \begin{tabular}{L{4.5cm}C{3cm}C{4cm}}
    \toprule
    \textbf{Nome e cognome} & \textbf{Matricola} & \textbf{Account GitHub} \\
    \midrule
    Alessandro Gremes  & 242330 & \code{alegrem123}      \\
    Paolo Sarcletti    & 246846 & \code{Puzzapol}        \\
    Alessandro Turri   & 244927 & \code{alessandroturri} \\
    \bottomrule
  \end{tabular}

\end{titlepage}

\hypersetup{pageanchor=true}
\pagenumbering{roman}
\tableofcontents
\clearpage
\pagenumbering{arabic}

% ==============================================================
\section{Sezione Introduttiva}
% ==============================================================

\subsection{Team Members}

\begin{table}[H]
  \centering\scriptsize
  \begin{tabular}{L{5cm}C{2.5cm}L{5cm}}
    \toprule
    \textbf{Nome e cognome} & \textbf{Matricola} & \textbf{Account GitHub} \\
    \midrule
    Alessandro Gremes  & 242330 & \code{alegrem123} \\
    Paolo Sarcletti    & 246846 & \code{Puzzapol} \\
    Alessandro Turri   & 244927 & \code{alessandroturri} \\
    \bottomrule
  \end{tabular}
\end{table}

\subsection{Project Idea}

RE-VALUE è una piattaforma web e mobile per favorire il riuso locale di oggetti ancora utilizzabili.
Un donatore pubblica un annuncio, un acquirente lo prenota e lo scambio fisico viene certificato tramite QR code.
Alla validazione il sistema completa atomicamente prenotazione e annuncio, aggiorna i wallet con crediti dinamici calcolati al momento della prenotazione, abilita recensioni, notifiche e reputazione.

\subsection{Link}

\begin{longtable}{L{4.2cm}L{10.5cm}}
\toprule
\textbf{Risorsa} & \textbf{Link / credenziali} \\
\midrule
\endhead
\bottomrule
\endfoot
Repository GitHub & \url{https://github.com/alegrem123/ReValue} \\
Link ad Apiary & \url{https://alessandrogremes.docs.apiary.io/#}; documentazione API storica del progetto, allineata alla specifica Apiary. \\
Specifica OpenAPI 3 / Swagger & \code{oas3.yaml} nel repository; fonte principale del contratto API machine-readable e importabile in Swagger Editor/SwaggerHub. \\
Swagger UI interattiva & \url{https://revalue-backend-84jb.onrender.com/api-docs/}; pagina generata da \code{oas3.yaml} per consultare e testare le API. \\
Backend in produzione & \url{https://revalue-backend-84jb.onrender.com} \\
Frontend in produzione & \url{https://revalue-frontend.onrender.com} \\
Account demo user & \code{demo.acquirente@revalue.local} / \code{Demo1234!} \\
Account demo admin & \code{demo.admin@revalue.local} / \code{Demo1234!} \\
\end{longtable}

Gli account demo sono generati tramite il comando \code{npm run seed:demo} nel backend. L'account utente standard dispone di un wallet iniziale di 160 crediti, cosi' da rendere dimostrabili anche il flusso premi e il riscatto dei coupon. L'account admin serve per accedere alla dashboard amministrativa web; nell'app mobile non abilita privilegi o schermate amministrative dedicate. Il seed crea inoltre altri utenti, annunci, prenotazioni e chat per rendere piu' popolato lo scenario dimostrativo.

% ==============================================================
\section{Sezione Generale}
% ==============================================================

\subsection{Strategia di branching}

\textbf{Sprint 1.} Il repository non ha usato una strategia \textit{master-only}: sono stati mantenuti due branch permanenti, \code{main} per la versione stabile e \code{preRelease} per l'integrazione. Alcuni commit sono comunque arrivati direttamente su \code{preRelease}; questa è stata una criticità di processo emersa nella retrospettiva e corretta nello Sprint 2.

\textbf{Sprint 2.} Il team ha reso obbligatori i branch temporanei per feature, bug fix o test prima dell'integrazione su \code{preRelease}. Il flusso usato è:
\begin{center}
\code{feat/fix/test/<nome>} $\rightarrow$ \code{preRelease} $\rightarrow$ \code{main}
\end{center}

I branch effettivamente presenti e usati a supporto della consegna sono:

\begin{itemize}[noitemsep]
  \item \code{main} --- branch stabile, sempre deployabile in produzione; aggiornato solo tramite merge da \code{preRelease} con pipeline CI verde.
  \item \code{preRelease} --- branch di integrazione; le feature vengono qui validate prima del rilascio su \code{main}.
  \item \code{feat/AT-*} --- branch settimanali di Alessandro Turri (settimane 4, 5, 6).
  \item \code{feat/ag-*} --- branch di Alessandro Gremes (admin, email, sicurezza, settimana 5).
  \item \code{feat/ps-*} --- branch di Paolo Sarcletti (recensioni, chat, frontend, Apiary).
  \item \code{fix/*} --- branch per bug fix urgenti, es.\ \code{fix/ag-priority-0}.
\end{itemize}

I messaggi di commit seguono in larga parte \textit{Conventional Commits} con scope per modulo, ad esempio \code{feat(premi): ...}, \code{fix(admin): ...}, \code{test(security): ...}, \code{docs(api): ...}. Alcuni commit di merge o di correzione veloce non sono perfettamente uniformi: sono stati mantenuti per preservare la cronologia reale del lavoro, ma la convenzione è stata usata come regola operativa.

Il CI/CD su GitHub Actions si attiva su push verso \code{main} e \code{preRelease}: esegue test backend con coverage, smoke test frontend, smoke test mobile e uno smoke test HTTP sul backend Render. Solo a pipeline verde si procede con il merge su \code{main}.

\subsection{Definizione di ``Done''}

La seguente definizione è stata adottata dal team a partire da Sprint 1 e rafforzata in Sprint 2 con i punti relativi a CI e deploy.

Una user story si considera \textbf{Done} quando:
\begin{enumerate}[noitemsep]
  \item Il contratto API è aggiornato in \code{oas3.yaml}; \code{apiary.apib} viene riallineato come versione API Blueprint pubblicabile su Apiary.
  \item Tutti i test automatici passano (\code{npm test} con 0 failing e copertura stabile).
  \item Nessuna regressione critica è stata introdotta nella suite di regressione esistente.
  \item[\textit{S2}] La pipeline CI GitHub Actions è verde prima del merge su \code{main}.
  \item[\textit{S2}] Il backend di produzione supera lo smoke test HTTP post-deploy.
\end{enumerate}

\subsection{Product Backlog}

Il product backlog completo è riportato nella tabella seguente. Le storie di Sprint 1 e Sprint 2 sono marchiate come \textit{Done}; le storie non ancora implementate rimangono nel backlog come lavoro futuro.

\begin{longtable}{L{0.9cm}L{6.5cm}C{1cm}C{1.3cm}C{1.3cm}L{1.8cm}}
\toprule
\textbf{ID} & \textbf{User Story} & \textbf{SP} & \textbf{Prio.} & \textbf{Sprint} & \textbf{Stato} \\
\midrule
\endhead
\bottomrule
\endfoot

% -- Sprint 1 --
PB-01 & Come utente, voglio registrarmi con email e password per accedere alla piattaforma. & 3 & Alta & S1 & Done \\
PB-02 & Come utente, voglio fare login e ricevere un token JWT per autenticarmi nelle richieste. & 3 & Alta & S1 & Done \\
PB-03 & Come donatore, voglio pubblicare un annuncio con oggetto, foto e data di scadenza. & 8 & Alta & S1 & Done \\
PB-04 & Come utente, voglio sfogliare il catalogo pubblico con filtri per categoria, dimensione e materiale. & 5 & Alta & S1 & Done \\
PB-05 & Come acquirente, voglio prenotare un oggetto disponibile. & 8 & Alta & S1 & Done \\
PB-06 & Come donatore, voglio generare un QR code per confermare il ritiro dell'oggetto. & 5 & Alta & S1 & Done \\
PB-07 & Come acquirente, voglio scansionare il QR code per completare lo scambio. & 5 & Alta & S1 & Done \\
PB-08 & Come utente, voglio vedere il mio saldo e lo storico transazioni del wallet. & 5 & Alta & S1 & Done \\
PB-09 & Come utente, voglio accedere alle funzionalità principali tramite il frontend web. & 8 & Alta & S1 & Done \\
PB-10 & Come utente, voglio accedere alle funzionalità principali tramite l'app mobile. & 5 & Media & S1 & Done \\
% -- Sprint 2 --
PB-11 & Come utente, voglio lasciare una recensione al partner dello scambio dopo il completamento. & 5 & Alta & S2 & Done \\
PB-12 & Come utente, voglio inviare messaggi al partner dello scambio tramite la chat integrata. & 8 & Alta & S2 & Done \\
PB-13 & Come utente, voglio ricevere notifiche in-app per eventi rilevanti (prenotazione, scambio, ecc.). & 5 & Alta & S2 & Done \\
PB-14 & Come sistema, voglio inviare push notification su dispositivo mobile tramite Expo. & 5 & Media & S2 & Parziale \\
PB-15 & Come utente, voglio riscattare coupon premio con i crediti accumulati. & 5 & Media & S2 & Done \\
PB-16 & Come utente, voglio segnalare comportamenti scorretti di altri utenti. & 3 & Media & S2 & Done \\
PB-17 & Come admin, voglio visualizzare le statistiche della piattaforma (scambi, crediti, utenti). & 5 & Alta & S2 & Done \\
PB-18 & Come admin, voglio sospendere, bannare in modo reversibile e riabilitare utenti. & 5 & Alta & S2 & Done \\
PB-19 & Come admin, voglio gestire e forzare lo stato degli annunci non conformi. & 3 & Media & S2 & Done \\
PB-20 & Come sistema, voglio applicare hardening sicurezza (Helmet, rate-limit, sanitizzazione). & 3 & Alta & S2 & Done \\
PB-21 & Come team, voglio CI/CD automatizzato su GitHub Actions con test e smoke produzione. & 3 & Alta & S2 & Done \\
PB-22 & Come team, voglio il deploy stabile su Render.com con MongoDB Atlas. & 5 & Alta & S2 & Done \\
PB-23 & Come utente, voglio poter aprire un ticket di supporto dalla piattaforma. & 2 & Bassa & S2 & Parziale API \\
% -- Backlog futuro --
PB-24 & Come utente, voglio accedere con SPID/SSO istituzionale. & 8 & Bassa & -- & Backlog \\
PB-25 & Come utente, voglio visualizzare gli annunci vicini su una mappa interattiva. & 5 & Bassa & -- & Backlog \\
PB-26 & Come utente, voglio ricevere notifiche in tempo reale via WebSocket (senza polling). & 5 & Bassa & -- & Backlog \\
PB-27 & Come utente mobile, voglio una schermata dedicata per i ticket di supporto. & 2 & Bassa & -- & Backlog \\
PB-28 & Come sistema, voglio verificare end-to-end la consegna push su dispositivo reale. & 3 & Media & -- & Backlog \\
PB-29 & Come sistema, voglio validare RNF2 con un load test documentato a 500 utenti concorrenti. & 5 & Media & -- & Backlog \\

\end{longtable}

% ==============================================================
\section{Sezione Sprint \#2}
% ==============================================================

\subsection{Goal}

\textbf{Periodo:} 18 maggio 2026 -- 7 giugno 2026.

Lo Sprint 2 ha avuto come obiettivo principale l'\textbf{estensione della piattaforma con le funzionalità sociali e di moderazione} e il \textbf{raggiungimento di un deploy stabile in produzione}. Concretamente: completare recensioni, chat, notifiche in-app, sistema premi e segnalazioni; realizzare il pannello di amministrazione con statistiche e moderazione utenti; irrobustire la sicurezza del backend; automatizzare test e deploy tramite CI/CD.

\subsection{Sprint 2 Planning}

\subsubsection*{Sprint 2 Backlog}

Il backlog di Sprint 2 corrisponde agli item PB-11---PB-23 del product backlog. La colonna ``owner'' indica il responsabile prevalente, mentre la colonna dei contributi riporta i casi in cui la feature è stata completata o corretta da più membri.

\begin{table}[H]
  \centering\scriptsize\sloppy
  \begin{tabular}{C{0.8cm}L{3.6cm}C{0.65cm}L{3cm}L{2.8cm}L{1.7cm}}
    \toprule
    \textbf{ID} & \textbf{User Story} & \textbf{SP} & \textbf{Owner prevalente / contributi} & \textbf{Stato} \\
    \midrule
    PB-11 & Recensioni post-scambio & 5 & P. Sarcletti; integrazione mobile AT & Done \\
    PB-12 & Chat e messaggi tra partecipanti & 8 & P. Sarcletti; notifiche AG; hardening AT & Done \\
    PB-13 & Notifiche in-app & 5 & A. Gremes; hook PS/AT & Done \\
    PB-14 & Push notification Expo & 5 & AG backend, AT mobile/test & Parziale \\
    PB-15 & Sistema premi / coupon & 5 & A. Turri & Done \\
    PB-16 & Segnalazioni utenti & 3 & A. Turri & Done \\
    PB-17 & Dashboard admin --- statistiche & 5 & A. Gremes; grafici/statistiche finali AT & Done \\
    PB-18 & Admin moderazione utenti & 5 & AG frontend iniziale; AT backend atomico/fix & Done \\
    PB-19 & Admin gestione annunci & 3 & A. Turri & Done \\
    PB-20 & Hardening sicurezza & 3 & AG hardening iniziale; AT regression/security fixes & Done \\
    PB-21 & CI/CD GitHub Actions & 3 & AG workflow iniziale; AT deploy/smoke Node 22 & Done \\
    PB-22 & Deploy Render.com + MongoDB Atlas & 5 & A. Turri & Done \\
    PB-23 & Ticket di supporto & 2 & A. Turri backend/API/test; UI dedicata futura & Parziale API \\
    \midrule
    \multicolumn{2}{r}{\textbf{Totale impegnato}} & \textbf{57} & & \multicolumn{2}{l}{53 SP Done pieni, 4 SP parziali dichiarati} \\
    \bottomrule
  \end{tabular}
\end{table}

\begin{itemize}[noitemsep]
  \item \textbf{Alessandro Gremes:} 38 commit nel periodo come \code{alegrem123} più 5 commit con autore \code{Alessandro Gremes}; aree principali: \code{/api/v1}, CORS, Apiary iniziale, email, notifiche, admin dashboard, CI, hardening errori/id.
  \item \textbf{Paolo Sarcletti:} 18 commit nel periodo; aree principali: recensioni backend, test recensioni, QR edge cases, chat search/image/polling/typing, frontend profilo pubblico e recensioni.
  \item \textbf{Alessandro Turri:} 143 commit nel periodo; aree principali: premi/coupon, segnalazioni/malus, mobile, deploy, admin backend, sicurezza finale, test regressione, documentazione D4 e tracciabilità.
\end{itemize}

Questa distribuzione spiega perché alcune user story sono attribuite a un owner prevalente ma riportano contributi incrociati: il lavoro finale è stato integrato e corretto collettivamente, soprattutto sui moduli admin, sicurezza, notifiche e mobile.

\subsubsection*{Burndown Chart}

Sprint 2 è durato 3 settimane (18 maggio -- 7 giugno 2026). La velocità iniziale è stata più bassa a causa della definizione dell'architettura dei moduli Social, poi ha accelerato nelle settimane centrali.

\begin{figure}[H]
\centering
\begin{tikzpicture}
\begin{axis}[
  width=13cm, height=7cm,
  xlabel={Giorno dello sprint},
  ylabel={Story point rimanenti},
  xmin=0, xmax=21,
  ymin=0, ymax=62,
  xtick={0,7,14,21},
  xticklabels={0 (16/5),7 (23/5),14 (30/5),21 (6/6)},
  ytick={0,10,20,30,40,50,57},
  grid=major,
  grid style={line width=0.3pt, draw=gray!30},
  legend pos=north east,
  legend style={font=\small},
  title={Burndown Sprint 2 -- RE-VALUE},
  title style={font=\normalsize\bfseries},
]
% Linea ideale
\addplot[color=gray, dashed, line width=1.5pt, mark=none]
  coordinates {(0,57)(21,0)};
\addlegendentry{Ideale}

% Linea effettiva
\addplot[color=blurevalue, solid, line width=2pt, mark=*, mark size=2.5pt]
  coordinates {(0,57)(7,49)(14,32)(21,4)};
\addlegendentry{Effettivo}

\end{axis}
\end{tikzpicture}
\caption{Burndown chart Sprint 2 (57 SP totali; 4 SP residui dichiarati: verifica push nativa su device reale e UI dedicata ticket supporto).}
\end{figure}

\subsection{Test Cases}

La verifica finale di Sprint 2 è stata eseguita il \textbf{8 giugno 2026} con i comandi:
\code{cd backend \&\& npm test -- --coverage --coverageReporters=text --runInBand},
\code{cd frontend \&\& npm test} e \code{cd mobile \&\& npm test}. L'esito complessivo è:
\textbf{backend 192/192 test passati su 21 suite}, \textbf{frontend 9/9 smoke test passati} e
\textbf{mobile 11/11 smoke test passati}. La coverage backend rilevata è
\textbf{69.48\% statements}, \textbf{57.39\% branches}, \textbf{67.63\% functions} e
\textbf{71.52\% lines}; i model e le route risultano quasi interamente coperti, mentre i controller
di chat/messaggi e gli adapter opzionali restano le aree con copertura più bassa.

\renewcommand{\arraystretch}{1.12}
\setlength\tabcolsep{2pt}
\hfuzz=10pt
\tiny
\begin{longtable}{L{1.2cm}C{0.45cm}L{2.35cm}L{2.05cm}L{1.95cm}L{0.9cm}L{2.3cm}L{1.55cm}L{2.05cm}}
\toprule
\textbf{User Story} & \textbf{N} & \textbf{Descrizione} & \textbf{Dati Test} & \textbf{Precondizioni} & \textbf{Dip.} & \textbf{Risultato Atteso} & \textbf{Ris. Eff.} & \textbf{Note} \\
\midrule
\endhead
\bottomrule
\endfoot

US01 & 1 & Registrazione valida. & Email valida, password 8+ char. & Email non registrata. & Auth & HTTP 201, JWT, wallet a 0, hash non esposto. & \code{PASS-auth} & EP; RF2/RF3, OCL \#1-\#2. \\
US01 & 2 & Email già registrata. & Stessa email usata due volte. & Utente esistente. & Auth & Registrazione rifiutata, nessun duplicato. & \code{PASS-auth} & EP; unicità email. \\
US01 & 3 & Password sotto soglia. & Password di 7 caratteri. & Email nuova. & Auth & HTTP 400, vincolo password applicato. & \code{PASS-auth} & BVA; RNF4/OCL \#2. \\
US01 & 4 & Login corretto. & Email/password valide. & Utente attivo. & Auth & HTTP 200 con JWT valido. & \code{PASS-auth} & EP; RF3. \\
US01 & 5 & Login con password errata. & Email valida, password errata. & Utente attivo. & Auth & HTTP 401, nessun token. & \code{PASS-auth} & EP; credenziali invalide. \\
US01 & 6 & Login utente bannato. & Credenziali corrette. & \code{bannato=true}. & Auth & HTTP 403, accesso negato. & \code{PASS-auth} & Cause-effect; moderazione. \\
US01 & 7 & JWT scaduto. & Token con \code{expiresIn=-1s}. & Route privata. & JWT & HTTP 401. & \code{PASS-auth} & Error guessing; sessioni. \\
US01 & 8 & Rate limit login. & 6 tentativi stesso identificatore. & Limiter attivo. & Auth & Il sesto tentativo riceve 429. & \code{PASS-auth} & BVA; soglia 5. \\

US02 & 9 & Catalogo pubblico. & GET \code{/api/v1/annunci}. & Annunci attivi nel DB. & Ann. & HTTP 200, lista paginata. & \code{PASS-annunci} & EP; RF4/UC8. \\
US02 & 10 & Privacy coordinate anonime. & Richiesta non autenticata. & Annuncio con lat/lng. & Ann. & Coordinate esatte non esposte. & \code{PASS-catalog} & Cause-effect; RF25. \\
US02 & 11 & Filtro categoria. & \code{?categoria=Mobili}. & Dataset misto. & Ann. & Solo annunci della categoria richiesta. & \code{PASS-annunci} & EP; RF22. \\
US02 & 12 & Filtro dimensione. & Dimensione \code{grande}. & Dataset misto. & Ann. & Solo annunci coerenti. & \code{PASS-annunci} & EP; RF22. \\
US02 & 13 & Filtro materiale. & \code{?materiale=legno}. & Dataset misto. & Ann. & Solo annunci coerenti. & \code{PASS-annunci} & EP; RF22. \\
US02 & 14 & Ordinamento per scadenza. & Tre scadenze diverse. & Annunci attivi. & DB & Ordine ASC per \code{dataScadenza}. & \code{PASS-annunci} & BVA/RNF1. \\

US03 & 15 & Creazione annuncio valida. & Titolo, oggetto, max 5 foto, scadenza futura. & Donatore autenticato. & Auth & HTTP 201, stato \code{DISPONIBILE}. & \code{PASS-annunci} & EP; RF15/RF16/OCL \#5. \\
US03 & 16 & Scadenza nel passato. & \code{dataScadenza < now}. & Donatore autenticato. & Auth & HTTP 400, annuncio non creato. & \code{PASS-annunci} & BVA; OCL \#5. \\
US03 & 17 & Foto oltre limite. & 6 foto base64. & Donatore autenticato. & Upload & HTTP 400. & \code{PASS-annunci} & BVA; max 5 foto. \\
US03 & 18 & Campi obbligatori assenti. & Titolo o categoria mancanti. & Donatore autenticato. & Model & HTTP 400. & \code{PASS-annunci} & EP; required fields. \\
US03 & 19 & Modifica annuncio disponibile. & PUT titolo/descrizione. & Donatore, stato \code{DISPONIBILE}. & Auth & HTTP 200, dati aggiornati. & \code{PASS-annunci} & EP; RF18/OCL \#8. \\
US03 & 20 & Modifica annuncio prenotato. & PUT su stato \code{PRENOTATO}. & Donatore proprietario. & Auth & HTTP 409. & \code{PASS-annunci} & Cause-effect; OCL \#8. \\
US03 & 21 & Soft delete. & DELETE annuncio disponibile. & Donatore proprietario. & Auth & \code{isAttivo=false}, documento conservato. & \code{PASS-annunci} & EP; soft delete. \\

US04 & 22 & Prenotazione valida. & Acquirente prenota annuncio disponibile. & Annuncio attivo e non scaduto. & Auth & HTTP 201, annuncio \code{PRENOTATO}. & \code{PASS-swap} & EP; OCL \#6-\#7. \\
US04 & 23 & Donatore prenota proprio annuncio. & ID annuncio proprio. & Donatore autenticato. & Auth & Prenotazione rifiutata. & \code{PASS-swap} & Cause-effect; OCL \#4. \\
US04 & 24 & Prenotazione concorrente. & Due acquirenti simultanei. & Annuncio disponibile. & DB & Un 201, un 409, una sola \code{ATTIVA}. & \code{PASS-conc} & Cause-effect; OCL \#9/RNF2. \\
US04 & 25 & Stress locale a 10 utenti. & Dieci prenotazioni simultanee. & Annuncio disponibile. & DB & Un successo, nove conflitti. & \code{PASS-conc} & Error guessing; concorrenza locale, RNF2 parziale. \\
US04 & 26 & Annullamento entro finestra. & DELETE prenotazione entro 15 min. & Prenotazione \code{ATTIVA}. & Auth & \code{ANNULLATA}, annuncio disponibile. & \code{PASS-pren} & BVA; OCL \#10-\#11. \\
US04 & 27 & Annullamento fuori finestra. & Prenotazione oltre 15 min. & Prenotazione \code{ATTIVA}. & Auth & HTTP 409 finestra scaduta. & \code{PASS-pren} & BVA; OCL \#10. \\

US05 & 28 & Generazione QR valida. & Donatore chiama \code{/qr/genera}. & Prenotazione \code{ATTIVA}. & Auth & Token QR non nullo, scadenza futura. & \code{PASS-swap} & EP; OCL \#13. \\
US05 & 29 & Acquirente tenta QR. & \code{/qr/genera} da acquirente. & Prenotazione attiva. & Auth & HTTP 403. & \code{PASS-qr} & Cause-effect; ruoli. \\
US05 & 30 & Validazione QR valida. & Codice QR corretto. & Prenotazione \code{ATTIVA}. & Wallet & \code{COMPLETATA}, annuncio \code{RITIRATO}, crediti dinamici accreditati a entrambi. & \code{PASS-swap} & Cause-effect; OCL \#14/RNF5; reverse auction. \\
US05 & 31 & QR già usato. & Stesso codice due volte. & Scambio completato. & QR & HTTP 409. & \code{PASS-qr} & Error guessing; idempotenza. \\
US05 & 32 & QR scaduto. & Token scaduto. & Documento QR esistente. & QR & HTTP 410. & \code{PASS-qr} & BVA; TTL 24h. \\
US05 & 33 & QR inesistente. & Codice casuale. & Nessun token. & QR & HTTP 404. & \code{PASS-qr} & EP; codice invalido. \\

US06 & 34 & Wallet nuovo. & Creazione utente. & Registrazione valida. & Auth & Bilancio 0, storico vuoto. & \code{PASS-wallet} & EP; OCL \#15. \\
US06 & 35 & Accredito positivo. & \code{ammontare=}. & Wallet esistente. & Wallet & Bilancio incrementato, transazione registrata. & \code{PASS-wallet} & EP; OCL \#16. \\
US06 & 36 & Accredito nullo. & \code{ammontare=0}. & Wallet esistente. & Wallet & Errore, saldo invariato. & \code{PASS-wallet} & BVA; OCL \#16. \\
US06 & 37 & Sottrazione valida. & Saldo 100, costo 50. & Wallet esistente. & Wallet & Saldo 50, transazione registrata. & \code{PASS-wallet} & EP; OCL \#17. \\
US06 & 38 & Saldo insufficiente. & Saldo 10, costo 50. & Wallet esistente. & Wallet & Errore, saldo non negativo. & \code{PASS-wallet} & BVA; OCL \#15/\#17. \\

US07 & 39 & Chat richiede JWT. & GET messaggi senza token. & Conversazione esistente. & Auth & HTTP 401. & \code{PASS-chat} & EP; RF14. \\
US07 & 40 & Utente estraneo legge chat. & GET da non partecipante. & Conversazione attiva. & Auth & HTTP 403. & \code{PASS-chat} & Cause-effect; RF13. \\
US07 & 41 & Partecipante legge storico. & GET da donatore/acquirente. & Prenotazione attiva. & Auth & HTTP 200 con messaggi. & \code{PASS-chat} & EP; RF11. \\
US07 & 42 & Partecipante invia messaggio. & POST testo valido. & Conversazione attiva. & Auth & HTTP 201, messaggio nello storico. & \code{PASS-chat} & EP; RF10. \\
US07 & 43 & Immagine oltre limite. & POST allegato grande. & Conversazione esistente. & Chat & HTTP 413. & \code{PASS-chat} & BVA; allegati. \\

US08 & 44 & Segnalazione valida. & Motivo non vuoto, target diverso. & Utente autenticato. & Auth & HTTP 201. & \code{PASS-segn} & EP; RF9/OCL \#18-\#19. \\
US08 & 45 & Auto-segnalazione. & Segnalante = segnalato. & Utente autenticato. & Auth & HTTP 409/400. & \code{PASS-segn} & Cause-effect; OCL \#19. \\
US08 & 46 & Motivo vuoto. & Motivo stringa vuota. & Utente autenticato. & Auth & HTTP 400. & \code{PASS-segn} & BVA; OCL \#18. \\
US08 & 47 & No-show da acquirente. & POST no-show da acquirente. & Prenotazione attiva. & Auth & HTTP 403. & \code{PASS-pren} & Cause-effect; RF19. \\
US08 & 48 & No-show da donatore. & POST no-show da donatore. & Prenotazione attiva. & Auth & Annullamento, annuncio disponibile, malus +1. & \code{PASS-pren} & EP; RF19/OCL \#20. \\
US08 & 49 & Terzo malus. & Applica terzo malus. & Utente con 2 malus. & Admin & Utente auto-sospeso. & \code{PASS-admin} & BVA; soglia 3/OCL \#20. \\

US09 & 50 & Lista premi attivi. & GET \code{/premi}. & Coupon attivi/non attivi. & Premi & Solo premi attivi ordinati per costo. & \code{PASS-premi} & EP; US09. \\
US09 & 51 & Riscatto saldo sufficiente. & Costo 50, saldo 100. & Utente autenticato. & Wallet & HTTP 201, codice univoco, stock aggiornato. & \code{PASS-premi} & EP; OCL \#17. \\
US09 & 52 & Riscatto saldo insufficiente. & Costo 50, saldo 10. & Utente autenticato. & Wallet & HTTP 409, saldo invariato. & \code{PASS-premi} & BVA; OCL \#17. \\
US09 & 53 & Riscatto già usato. & PATCH usato due volte. & Riscatto proprietario. & Premi & Secondo update rifiutato. & \code{PASS-premi} & Error guessing; stato. \\

US10 & 54 & Utente normale accede admin. & JWT ruolo \code{user}. & Route admin. & Auth & HTTP 403. & \code{PASS-security} & EP; ruoli. \\
US10 & 55 & Sospensione e riabilitazione. & Admin sospende poi riabilita. & Utente non bannato. & Admin & Stato aggiornato, accessi coerenti. & \code{PASS-admin} & EP; moderazione. \\
US10 & 56 & Riabilitazione bannato. & Admin riabilita utente bannato. & \code{bannato=true}. & Admin & Blocco rimosso e login nuovamente consentito. & \code{PASS-admin} & Cause-effect; decisione S2: ban reversibile. \\
US10 & 57 & Forzatura stato annuncio. & Admin forza annuncio bloccato. & Prenotazione attiva. & Admin & Stato aggiornato e prenotazioni annullate. & \code{PASS-admin} & Error guessing; RF31. \\

US11 & 58 & Recensione prima dello scambio. & POST su prenotazione attiva. & Scambio non completato. & Auth & HTTP 409. & \code{PASS-review} & Cause-effect; OCL \#21. \\
US11 & 59 & Recensione post-scambio. & POST dopo QR valido. & Prenotazione completata. & Auth & HTTP 201. & \code{PASS-review} & EP; RF21/RF28. \\
US11 & 60 & Recensione duplicata. & Stesso autore/prenotazione. & Prima recensione presente. & Auth & HTTP 409. & \code{PASS-review} & Error guessing; duplicati. \\

US12 & 61 & Lista notifiche richiede token. & GET notifiche anonimo. & Notifiche esistenti. & Auth & HTTP 401. & \code{PASS-notif} & EP; RF12. \\
US12 & 62 & Lista notifiche filtrabile. & \code{?letta=false}. & Notifiche miste. & Notif. & HTTP 200, paginazione coerente. & \code{PASS-notif} & EP; RF12. \\
US12 & 63 & Marca singola letta. & PATCH notifica propria. & Notifica non letta. & Auth & Solo notifica propria aggiornata. & \code{PASS-notif} & Cause-effect; ownership. \\
US12 & 64 & Token Expo valido. & Token Expo valido. & Utente autenticato. & Mobile & Token salvato nel profilo. & \code{PASS-push} & EP; push parziale. \\
US12 & 65 & Token Expo non valido. & Stringa qualsiasi. & Utente autenticato. & Mobile & HTTP 400. & \code{PASS-push} & EP; validazione token. \\

US13 & 66 & Ticket supporto valido. & Oggetto e messaggio non vuoti. & Utente autenticato. & Supp. & HTTP 201. & \code{PASS-supp} & EP; supporto S2. \\
US13 & 67 & Ticket supporto vuoto. & Messaggio vuoto. & Utente autenticato. & Supp. & HTTP 400. & \code{PASS-supp} & BVA; required. \\
US13 & 68 & Lista ticket personali. & GET ticket utente. & Ticket di più utenti. & Supp. & Solo ticket del richiedente. & \code{PASS-supp} & Cause-effect; privacy. \\

S2-N1 & 69 & Endpoint legacy scambi. & GET/POST \code{/api/v1/scambi/*}. & Route legacy. & API & HTTP 410 con \code{Deprecation}. & \code{PASS-legacy} & Error guessing; compatibilità. \\
S2-N2 & 70 & Smoke frontend. & Node test su viste web. & Repository completo. & Web & 8 smoke test passati. & \code{PASS-front} & Regression; frontend. \\
S2-N3 & 71 & Smoke mobile. & Node test su schermate RN. & Repository completo. & Mobile & 10 smoke test passati. & \code{PASS-mobile} & Regression; mobile. \\
S2-N4 & 72 & CI backend con coverage. & GitHub Actions su push/PR. & Branch \code{main/preRelease}. & CI & Test backend, frontend e mobile eseguiti. & \code{PASS-CI} & White-box; regressione. \\
S2-N5 & 73 & Push end-to-end device reale. & Device fisico Expo/EAS. & Permessi e build configurati. & Expo & Notifica ricevuta sul device. & \code{PENDING-PB28} & Manuale; dipende da EAS. \\

\end{longtable}
\setlength\tabcolsep{6pt}
\renewcommand{\arraystretch}{1}
\normalsize

\subsection{Sprint Review}

La Sprint Review si è svolta al termine della terza settimana di Sprint 2 (8 giugno 2026). Durante la demo è stato mostrato l'intero flusso della piattaforma in produzione: registrazione, login, pubblicazione annuncio, prenotazione, generazione QR, scansione, completamento scambio con aggiornamento wallet tramite crediti dinamici, invio messaggio in chat, riscatto coupon, invio segnalazione e gestione da pannello admin (statistiche, moderazione utente, forzatura stato annuncio).

\medskip
\noindent\textbf{Aspetti rilevanti della demo:}

\begin{itemize}[noitemsep]
  \item Tutti i flussi principali sono risultati operativi in produzione; la verifica del 8 giugno ha confermato backend 192/192, frontend 9/9 e mobile 11/11 test passati.
  \item Il pannello admin ha mostrato le statistiche in tempo reale calcolate dalla suite di aggregazioni MongoDB.
  \item La demo della chat ha evidenziato il funzionamento del middleware \code{requireParticipant}: utenti estranei non possono accedere allo storico.
  \item Le notifiche in-app vengono create e segnate come lette correttamente; il canale push su dispositivo reale non è stato dimostrato live per dipendenza dalla configurazione EAS.
\end{itemize}

\medskip
\noindent\textbf{Discussione:}
Nel complesso la demo ha soddisfatto i criteri di accettazione dei principali use case. Le limitazioni emerse e dichiarate nel report sono due: la consegna push nativa non è stata verificata end-to-end su dispositivo fisico con build EAS, e il modulo supporto è completo lato backend/API ma non ha ancora una schermata dedicata comparabile alle altre aree web/mobile. Sono stati predisposti due account demo verificati in produzione per rendere la valutazione ripetibile senza interventi manuali su MongoDB Atlas.

\subsection{Product Backlog Refinement}

Il meeting di refinement si è tenuto al termine della seconda settimana di Sprint 2. Le principali variazioni al product backlog sono state le seguenti.

\begin{itemize}[noitemsep]
  \item Le storie relative a SPID/SSO e mappa OpenStreetMap (PB-24, PB-25) sono state spostate in icebox: il rapporto tra effort e valore prototipale non giustifica l'integrazione entro la deadline D4.
  \item La storia PB-14 (push notification nativa su device reale) è stata ridimensionata: il backend Expo Push e il package \code{expo-notifications} nel client mobile sono presenti; la verifica end-to-end su dispositivo reale (PB-28) è stata separata e spostata nel backlog futuro.
  \item È stata aggiunta la storia PB-23 (ticket di supporto) dopo il feedback ricevuto durante la Sprint Review intermedia. Entro Sprint 2 sono stati completati backend, API e test; la UI dedicata è stata spostata nel backlog futuro per non ridurre la qualità dei flussi core.
  \item Il sistema crediti è stato evoluto da accredito fisso a \textit{reverse auction}: i punti assegnati a donatore e acquirente sono calcolati dinamicamente in base a categoria e vicinanza alla scadenza, poi salvati sulla prenotazione e accreditati solo alla validazione QR.
  \item La moderazione PB-18 è stata interpretata come blocco amministrativo reversibile: l'utente bannato non può accedere finché il flag \code{bannato=true}, ma un admin può riabilitarlo in caso di errore o revisione manuale. Questa scelta mantiene tracciabilità e recuperabilità dell'account nel prototipo accademico.
  \item È stato identificato un bug nella logica dello scheduler (annunci scaduti non venivano marcati \code{isAttivo=false}) e nel metodo \code{forzaStatoAnnuncio} (non ripristinava \code{isAttivo=true}): entrambi i fix sono stati inclusi nello sprint corrente come task tecnici.
\end{itemize}

\subsection{Sprint Retrospective}

La retrospettiva si è tenuta l' 8 giugno 2026, a valle della Sprint Review e della preparazione della documentazione D4.

\medskip
\noindent\textbf{Cosa ha funzionato:}
\begin{itemize}[noitemsep]
  \item La separazione in layer (route → controller → service → model) ha permesso di aggiungere tutti i moduli Sprint 2 senza riscrivere il backend: ogni membro ha lavorato su file distinti con conflitti minimi.
  \item L'adozione di Conventional Commits e branch per membro ha reso la cronologia git leggibile e ha semplificato la costruzione della lista commit finale.
  \item La suite di test automatici (backend 192/192, frontend 9/9, mobile 11/11) ha dato fiducia nelle modifiche: ogni merge su \code{preRelease} era verificato dal CI prima di toccare \code{main}.
  \item Il deploy su Render.com ha richiesto meno effort del previsto; la configurazione con variabili d'ambiente ha tenuto le credenziali fuori dal repository.
\end{itemize}

\medskip
\noindent\textbf{Cosa non ha funzionato:}
\begin{itemize}[noitemsep]
  \item La gestione del time-boxing non è stata sempre rispettata: alcune storie (chat, admin moderazione) si sono rivelate più complesse del previsto e hanno richiesto più SP di quanto stimato.
  \item L'integrazione push nativa ha richiesto tempo su configurazione EAS che non era prevista inizialmente; il risultato è parziale.
  \item I bug nello scheduler e in \code{forzaStatoAnnuncio} sono stati scoperti tardi: un code review più attento sulle funzioni critiche avrebbe permesso di intercettarli prima.
\end{itemize}

\medskip
\noindent\textbf{Adozione di pratiche agili:}
  Durante lo sviluppo di RE-VALUE il team ha adottato un processo agile ispirato a Scrum, adattato al contesto accademico e alla dimensione ridotta del gruppo. Il lavoro è stato organizzato in sprint
  incrementali, con un backlog mantenuto e aggiornato in base alle priorità funzionali, ai vincoli del corso.

  All'inizio dello Sprint 2 il team ha svolto una fase di planning per selezionare le user story più rilevanti: chat, notifiche, recensioni, premi, segnalazioni, moderazione admin, deploy e hardening API. A
  ciascuna attività sono stati associati story point e un responsabile prevalente, pur mantenendo collaborazione trasversale nei punti più critici, come QR, wallet, sicurezza e documentazione.

  La gestione del lavoro è stata supportata da una strategia Gitflow semplificata: \code{main} come branch stabile, \code{preRelease} come branch di integrazione e branch temporanei \code{feat/*}, \code{fix/
  *} e \code{test/*} per lo sviluppo delle singole attività. Questa scelta ha permesso di integrare progressivamente le funzionalità e di mantenere una cronologia leggibile tramite Conventional Commits.

  Il team ha applicato una Definition of Done condivisa: una user story è stata considerata conclusa solo dopo aggiornamento del contratto API, implementazione, test automatici o smoke test, verifica su
  ambiente integrato e assenza di regressioni critiche. Per le funzionalità backend più rischiose sono stati aggiunti test di regressione con Jest/Supertest, mentre frontend e mobile sono stati verificati
  con smoke test automatici e prove manuali sui flussi principali.

  Durante lo sprint sono stati svolti momenti frequenti di sincronizzazione informale, usati per riallineare priorità, risolvere blocchi tecnici e ridistribuire attività quando alcune stime iniziali si sono
  rivelate ottimistiche. La retrospettiva di Sprint 1 ha portato a decisioni concrete in Sprint 2: introduzione sistematica del prefisso \code{/api/v1}, maggiore attenzione alla documentazione API,
  rafforzamento dei test su concorrenza e sicurezza, e deploy stabile su cloud.

  Il processo non è stato applicato in modo puramente formale: alcune revisioni sono state svolte direttamente tramite integrazione su \code{preRelease} e controllo con CI, invece che sempre tramite pull
  request strutturate. Questa deviazione è stata documentata come adattamento pragmatico alla dimensione del team, mantenendo comunque i principi agili di incremento funzionante, ispezione continua,
  adattamento e trasparenza sulle decisioni tecniche.


\medskip
\noindent\textbf{Spunti per eventuali sprint futuri:}
\begin{itemize}[noitemsep]
  \item Completare la verifica push end-to-end su dispositivo fisico con EAS.
  \item Eseguire e documentare un load test RNF2 con 500 utenti concorrenti simulati.
  \item Valutare WebSocket per le notifiche in tempo reale (attualmente il frontend fa polling).
  \item Aggiungere indice geospaziale 2dsphere per ottimizzare le query di catalogo con filtro distanza.
  \item Migrare l'hashing password da SHA-256 a bcrypt/Argon2 in un contesto di produzione reale.
\end{itemize}

% ==============================================================
\section{Sezione Finale}
% ==============================================================

\subsection{Diagramma del deploy finale}
Il sistema è deployato su un'architettura cloud a tre blocchi principali: repository e pipeline CI/CD su GitHub, servizi applicativi su Render e persistenza su MongoDB Atlas. Il branch \code{main}
  rappresenta la versione stabile consegnabile: dopo il merge da \code{preRelease}, GitHub Actions esegue test backend con coverage, smoke test frontend/mobile e controlli HTTP sugli ambienti di produzione.
  Solo dopo una pipeline verde il deploy Render viene considerato valido per la consegna.

  \begin{figure}[H]
    \centering
    \includegraphics[width=\textwidth]{diagram.png}
    \caption{Deployment diagram di RE-VALUE: GitHub Actions, servizi Render, client web/mobile e MongoDB Atlas.}
    \label{fig:deployment}
  \end{figure}

  Il backend è eseguito come Web Service Render e contiene l'API Node.js/Express esposta sotto il prefisso \code{/api/v1}, la Swagger UI accessibile da \code{/api-docs} e lo scheduler batch per la scadenza
  automatica degli annunci. La configurazione sensibile è gestita tramite variabili d'ambiente Render, tra cui \code{MONGODB\_URI}, \code{JWT\_SECRET}, \code{FRONTEND\_URL}, \code{NODE\_ENV=production} e
  \code{REQUEST\_BODY\_LIMIT}.

  Il frontend web è deployato come sito statico Render e comunica con il backend tramite chiamate HTTPS REST verso gli endpoint versionati. L'app mobile Expo/React Native usa lo stesso contratto API e si
  collega al backend di produzione tramite HTTPS. La persistenza è affidata a MongoDB Atlas, raggiunto dal backend tramite Mongoose e connection string TLS. Le principali collezioni includono utenti,
  annunci, prenotazioni, wallet, token QR, conversazioni, recensioni, segnalazioni, notifiche, coupon e riscatti.

  Questa configurazione separa chiaramente client, API e database, mantiene le credenziali fuori dal repository e rende verificabile il rilascio tramite CI/CD e smoke test su produzione.


\subsection{Note sui servizi esterni}

\begin{itemize}[noitemsep]
  \item \textbf{SPID/SSO:} non implementato nel prototipo. L'autenticazione usa email/password con JWT. Dichiarato come scelta progettuale per il contesto accademico.
  \item \textbf{OpenStreetMap (backend):} non attivo come servizio backend. La geolocalizzazione rimane lato client; le coordinate degli annunci vengono filtrate via haversine in-memory per il catalogo con raggio.
  \item \textbf{Email (Nodemailer/SMTP):} adapter opzionale. Invia solo se \code{SMTP\_USER} e \code{SMTP\_PASS} sono configurati; salta senza bloccare il flusso applicativo in assenza di configurazione.
  \item \textbf{Expo Push API:} il backend chiama \code{https://exp.host/--/api/v2/push/send} in modalità fire-and-forget quando l'utente ha un \code{expoPushToken} valido. Il package \code{expo-notifications} è presente nel client mobile; la verifica end-to-end su dispositivo fisico con configurazione EAS è stata rimandata al backlog futuro (PB-28).
  \item \textbf{SHA-256 salato per hashing password:} scelta deliberata per il prototipo accademico, in coerenza con il vincolo OCL \#2 implementato nel modello utente: il backend accetta hash legacy a 64 caratteri o il formato salato \code{salt:sha256}. In un contesto produttivo si raccomanda bcrypt o Argon2 con parametri di costo configurabili.
\end{itemize}

\subsection{Stack tecnologico}

\subsubsection*{Backend}

\begin{itemize}[noitemsep]
  \item \textbf{Runtime:} Node.js 22 (LTS) --- versione usata in ambiente CI/CD; lo sviluppo locale è stato condotto su versioni recenti compatibili.
  \item \textbf{Framework:} Express 5.2
  \item \textbf{Database:} MongoDB 7, Mongoose 8 (ODM)
  \item \textbf{Autenticazione:} JWT (\code{jsonwebtoken} 9), hashing SHA-256 (\code{crypto} nativo)
  \item \textbf{Sicurezza:} Helmet 8, express-rate-limit 8, express-mongo-sanitize 2, CORS configurato per ambiente
  \item \textbf{Email:} Nodemailer 8 (adapter opzionale)
  \item \textbf{Logging:} Morgan 1
  \item \textbf{Testing:} Jest 29, Supertest 7, mongodb-memory-server 9 (MongoMemoryReplSet per transazioni)
  \item \textbf{CI:} GitHub Actions (Node.js 22, coverage obbligatoria)
  \item \textbf{Deploy:} Render.com Web Service
\end{itemize}

\subsubsection*{Frontend web}

\begin{itemize}[noitemsep]
  \item HTML5, CSS3, JavaScript ES modules (vanilla, nessun framework)
  \item Bootstrap (layout e componenti UI)
  \item Chart.js (grafici nella dashboard admin)
  \item \code{jsQR} (decodifica QR da browser)
  \item Client API centralizzato (\code{frontend/js/apiClient.js})
  \item Deploy: Render.com Static Site
\end{itemize}

\subsubsection*{App mobile}

\begin{itemize}[noitemsep]
  \item Expo 54, React Native 0.81.5
  \item \code{@react-native-async-storage/async-storage} (persistenza token)
  \item \code{expo-camera 17} (scansione QR)
  \item \code{expo-image-picker 17} (foto annunci)
  \item \code{expo-location 19} (geolocalizzazione)
  \item \code{expo-notifications 0.32} (push notification, installato, integrazione end-to-end su device reale in PB-28)
  \item \code{react-native-qrcode-svg 6}, \code{react-native-svg 15} (generazione QR)
  \item \code{@expo/vector-icons 15}
\end{itemize}

\subsubsection*{Documentazione e strumenti}

\begin{itemize}[noitemsep]
  \item OpenAPI 3 / Swagger (\code{oas3.yaml}) --- fonte formale e machine-readable del contratto API
  \item API Blueprint / Apiary (\code{apiary.apib}) --- documentazione pubblicabile su Apiary
  \item Conventional Commits per la cronologia git
  \item GitHub Actions per CI/CD con smoke test su produzione
\end{itemize}


\subsection{Conclusioni}

RE-VALUE è partito come un nucleo minimale (auth, annunci, prenotazione, QR, wallet) e nel corso dei due sprint si è evoluto in una piattaforma completa con funzionalità sociali, moderazione admin, sistema reputazionale, crediti dinamici e deploy in produzione su cloud. La suite automatica finale passa con 191 test backend su 21 suite, 9 smoke test frontend e 11 smoke test mobile, garantendo un livello di confidenza adeguato per un prototipo accademico.

Le scelte tecnologiche più discusse --- SHA-256 per hashing, assenza di SPID/SSO, push nativa non verificata end-to-end su device reale --- sono state documentate esplicitamente come decisioni deliberate per il contesto di prototipo, con indicazione delle evoluzioni consigliate per un contesto produttivo.

Dal punto di vista del processo, l'adozione di Gitflow semplificato e CI/CD obbligatoria ha nettamente migliorato la qualità delle integrazioni rispetto a Sprint 1. La comunicazione interna è stata efficace, anche se le stime iniziali degli SP per alcuni moduli (chat, admin) si sono rivelate ottimistiche.

Il team considera gli obiettivi di Sprint 2 sostanzialmente raggiunti: tutte le user story ad alta priorità sono Done, i flussi principali sono verificabili in produzione, la documentazione Apiary/OpenAPI è allineata al backend e il report integra test cases e matrice di tracciabilità.

% ==============================================================
\appendix
\section{Matrice di tracciabilità requisiti}
% ==============================================================

La tabella seguente mappa ogni requisito funzionale (RF), use case (UC) e vincolo OCL ai moduli backend, alle schermate frontend/mobile e alle evidenze di test. Gli stati sono: \textbf{Completo}, \textbf{Parziale}, \textbf{Prototipale}, \textbf{Futuro}.

\newcommand{\tr}[1]{{\ttfamily\seqsplit{#1}}}

\setlength\tabcolsep{3pt}
\renewcommand{\arraystretch}{1.15}
\scriptsize
\begin{longtable}{L{2.0cm}L{2.2cm}L{3.4cm}L{2.5cm}L{2.2cm}L{1.4cm}}
\toprule
\textbf{Area} & \textbf{RF/UC/OCL} & \textbf{Backend / API} & \textbf{Frontend / Mobile} & \textbf{Test} & \textbf{Stato} \\
\midrule
\endhead
\bottomrule
\endfoot

Autenticazione
  & Auth
  & \tr{authRoutes, authController,} \tr{jwt.js, authService}
  & \tr{login.html, register.html,} \tr{AuthScreen}
  & \tr{auth.test.js,} \tr{security.test.js}
  & Completo \\

Password hash
  & OCL \#2 / RNF
  & \tr{password.js, userModel.js}
  & login/register web e mobile
  & \tr{auth.test.js}
  & Prototipale\footnotemark[1] \\

Profilo utente
  & Users
  & \tr{usersRoutes, usersController}
  & \tr{profile.html, ProfileScreen}
  & smoke
  & Completo \\

Profilo pubblico
  & RF8
  & \tr{GET /users/:id/profilo,} \tr{recensioniController}
  & \tr{public-profile.html}
  & \tr{reviews.test.js}
  & Completo \\

Catalogo
  & RF4 / UC8 / RF22
  & \tr{GET /annunci,} \tr{annunciController}
  & \tr{catalog.html, CatalogScreen}
  & \tr{annunci.test.js,} \tr{catalog\_flow.test.js}
  & Completo \\

Crea annuncio
  & RF15 / RF16 / OCL \#5
  & \tr{POST /annunci,} \tr{annuncioModel}
  & \tr{create-annuncio.html,} \tr{CreateAnnuncioScreen}
  & \tr{annunci.test.js,} \tr{catalog\_flow.test.js}
  & Completo \\

Modifica annuncio
  & RF18 / OCL \#7
  & \tr{PUT /annunci/:id}
  & \tr{edit-annuncio.html,} \tr{MyAnnunciScreen}
  & \tr{annunci.test.js}
  & Completo \\

Cancellazione
  & RF18 / OCL \#8
  & \tr{DELETE /annunci/:id}
  & \tr{my-annunci.html}
  & \tr{annunci.test.js,} \tr{catalog\_flow.test.js}
  & Completo \\

Stato annuncio
  & state machine
  & \tr{PATCH /annunci/:id/stato}
  & viste donatore
  & \tr{annunci.test.js}
  & Completo \\

Prenotazione
  & UC2 / OCL \#4,\#6, \#9,\#10
  & \tr{prenotazioniRoutes,} \tr{prenotazioniController,} \tr{prenotazioniService}
  & \tr{annuncio.html,} \tr{MyBookingsScreen}
  & \tr{prenotazioni.test.js,} \tr{concurrent.test.js}
  & Completo \\

Concorrenza
  & OCL \#9,\#10 / RNF2 parziale
  & update atomico su campo \tr{versione}
  & n/a
  & \tr{concurrent.test.js}
  & Parziale \\

Annullamento
  & RF19 / RF20 / OCL \#11
  & disdetta e no-show routes
  & \tr{mybookings.html,} \tr{MyBookingsScreen}
  & \tr{prenotazioni.test.js,} \tr{prenotazioni\_flow.test.js}
  & Completo \\

QR scambio
  & UC3 / RF17,RF27 / OCL \#13--\#15
  & \tr{qrRoutes,} \tr{scambioQrService}
  & \tr{qr-display.html,} \tr{qr-scan.html,} \tr{QRScanScreen}
  & \tr{swap.test.js,} \tr{qr\_edge\_cases.test.js}
  & Completo \\

Wallet
  & RF5,RF6 / OCL \#16,\#17
  & \tr{walletRoutes, walletService}
  & \tr{wallethistory.html}
  & \tr{wallet.test.js,} \tr{swap.test.js}
  & Completo \\

Premi / coupon
  & UC7 / OCL \#17
  & \tr{premiController,} \tr{couponModel, riscattoModel}
  & \tr{premi.html, my-premi.html,} \tr{PremiScreen}
  & \tr{premi.test.js}
  & Completo \\

Recensioni
  & RF8,RF21,RF28 / OCL \#21
  & \tr{recensioniController}
  & \tr{swap-success.html,} \tr{public-profile.html,} \tr{SwapSuccessScreen}
  & \tr{reviews.test.js}
  & Completo \\

Chat
  & RF10,RF11, RF13,RF14
  & \tr{chatRoutes, messaggiRoutes,} \tr{requireParticipant}
  & \tr{chat.html, messaggi.html,} \tr{ChatScreen}
  & \tr{chat\_access.test.js,} smoke
  & Completo \\

Notifiche in-app
  & RF12
  & \tr{notificheService,} \tr{notificaModel}
  & \tr{notifiche.html,} \tr{NotificheScreen}
  & \tr{notifiche.test.js}
  & Completo \\

Push native
  & RF12 esteso
  & \tr{sendExpoPushIfConfigured,} \tr{PATCH /users/me/push-token}
  & token mobile (parziale)
  & \tr{support\_push.test.js}
  & Parziale\footnotemark[2] \\

Email
  & servizio esterno
  & \tr{emailService.js}
  & n/a
  & doc servizi
  & Prototipale \\

Segnalazioni
  & RF9 / OCL \#18,\#19
  & \tr{segnalazioniController}
  & \tr{segnala.html, my-reports.html}
  & \tr{segnalazioni.test.js,} \tr{security.test.js}
  & Completo \\

Moderazione admin
  & RF29 / OCL \#20
  & \tr{adminController, userService}
  & \tr{admin/dashboard.html}
  & \tr{admin\_moderation.test.js}
  & Completo \\

Dashboard admin
  & RF30,RF31 / UC14
  & \tr{GET /admin/statistiche}
  & dashboard admin
  & \tr{admin\_moderation.test.js}
  & Completo \\

Ticket supporto
  & S2
  & \tr{supportoController,} \tr{ticketSupportoModel}
  & nessuna UI dedicata
  & \tr{support\_push.test.js}
  & Parziale\footnotemark[3] \\

Contratto API
  & REST / Apiary
  & \tr{apiary.apib, /api/v1/*}
  & client web/mobile
  & smoke API client
  & Completo \\

App mobile
  & prototipo
  & \tr{App.js,} \tr{mobile/src/screens/}
  & React Native/Expo
  & \tr{mobile/tests/smoke}
  & Completo prototipo \\

CI/CD
  & processo
  & \tr{.github/workflows/ci.yml}
  & deploy Render
  & GitHub Actions
  & Completo \\

SPID/SSO, OSM
  & feature future
  & n/a
  & n/a
  & n/a
  & Futuro \\

\end{longtable}
\setlength\tabcolsep{6pt}
\renewcommand{\arraystretch}{1}
\normalsize

\footnotetext[1]{SHA-256 è scelta accademica dichiarata; bcrypt/Argon2 consigliati per produzione.}
\footnotetext[2]{Backend pronto e testato; verifica end-to-end su device fisico (EAS) rinviata a PB-28.}
\footnotetext[3]{Backend implementato e testato; schermata UI mobile dedicata in PB-27.}

\end{document}
